% TOP_PROBLEM: {"id": 3258, "title": "Seymour's Second Neighbourhood Conjecture", "category_id": 3, "category": {"id": 3, "name": "graph_theory", "display_name": "Graph Theory", "description": "Problems involving graphs, networks, and their properties.", "slug": "graph-theory", "order_index": 3, "created_at": "2026-07-31T15:26:25.671Z"}, "status": "open", "problem_number": "OPG-646", "set_id": 12}
% FIRST_POSTED: 2026-10-01
% ENTRY_KIND: statement_only
% UnsolvedMath ID 3258; source catalogue code OPG-646
% !TeX program = lualatex
% Definition and sources only; this file is not an LLM solution attempt.
\documentclass[11pt,a4paper]{article}
\usepackage[margin=25mm]{geometry}
\usepackage{fontspec}
\setmainfont{lmroman10-regular.otf}[BoldFont=lmroman10-bold.otf,
 ItalicFont=lmroman10-italic.otf,BoldItalicFont=lmroman10-bolditalic.otf]
\usepackage{amsmath,amssymb,mathtools}
\usepackage{unicode-math}
\setmathfont{latinmodern-math.otf}
\AtBeginDocument{\let\setminus\smallsetminus}
\usepackage{xurl,hyperref}
\hypersetup{colorlinks=true,urlcolor=blue}
\setcounter{secnumdepth}{0}
\setlength{\parindent}{0pt}
\setlength{\parskip}{5pt}
\setlength{\emergencystretch}{3em}
\begin{document}
\section*{Seymour's Second Neighbourhood Conjecture}\label{3258}
{\small UnsolvedMath ID 3258 · OPG-646 · Graph Theory · OpenGarden}
\subsection{Definitions and mathematical statement}
Let $D=(V,A)$ be a finite nonempty oriented graph: it has no loops, no
parallel arcs, and never contains both $uv$ and $vu$. The first
outneighbourhood of $v$ is
$N^+(v)=\{u\in V:vu\in A\}$. Define its second outneighbourhood by
\[
 N^{++}(v)=\{w\in V\setminus(\{v\}\cup N^+(v)):
                  \exists u\in V\text{ with }vu,uw\in A\}.
\]
Thus $N^{++}(v)$ consists of vertices at directed distance exactly two
from $v$, not all vertices reachable by a two-step walk. A vertex already
in $N^+(v)$ is excluded even if it also has such a walk. Both sets count
distinct vertices, without multiplicity.

Seymour's second neighbourhood conjecture asks whether every finite
nonempty oriented graph has at least one vertex $v$ satisfying
\[
                       |N^+(v)|\le |N^{++}(v)|.
\]
The desired vertex may depend on the entire orientation. There is no
claim that every vertex, or a prescribed vertex, satisfies the inequality.
A sink has no outneighbours and therefore already satisfies it with
both sides zero.

The orientation condition cannot simply be dropped: in the complete
symmetric digraph every other vertex is a first outneighbour, leaving
no second outneighbours in the exact-distance sense. The asserted bound
has coefficient exactly one; a positive lower ratio strictly below one
would be a different, weaker statement. All finite orders and all
underlying simple graphs are included.
\subsection{Short English statement}
Must every oriented graph contain a vertex with at least as many vertices at outward distance two as at outward distance one?
\subsection{Sources}
\begin{itemize}
\item[S1] ULAM.AI and the UnsolvedMath Contributors, supplied problems.json, record 3258 (OPG-646), OpenGarden collection. Catalogue identity and source transcription; CC BY 4.0.
\url{https://huggingface.co/datasets/ulamai/UnsolvedMath}.
\item[S2] Open Problem Garden, Seymour's Second Neighbourhood Conjecture. Original problem and discussion; the mathematical formulation below is expanded from the supplied source record.
\url{http://www.openproblemgarden.org/op/seymours_second_neighbourhood_conjecture}.
\end{itemize}
\end{document}
